\section{What an independently positive realization requires}\label{sec:positive}
\subsection{State norms, constructive hypotheses, and linear sources}
\begin{definition}[Exact positive realization]\label{def:realization}
An exact positive realization of the Weil form consists of a specified
model $M$ with a positive Hilbert space $\HH_M$ and complex-linear maps
$\Phi_{M,L}:C_c^\infty(I_L)\to\HH_M$ such that
\begin{equation}\label{eq:exact-match}
 \ip{\Phi_{M,L}(f)}{\Phi_{M,L}(g)}=Q_L(f,g)
 \qquad (L>0).
\end{equation}
The model, its physical inner product, source rule and normalization
must be specified independently of an assumed positive Weil form.
Construction of the model may be an explicit hypothesis; the identity
\eqref{eq:exact-match} must be derived from the specified data.
\end{definition}

One usual form is $\Phi_{M,L}(f)=\mathcal O_{M,L}(f)\Omega_M$.
The observable $\mathcal O_{M,L}(f)$ need not be a positive operator.
The positivity argument is
\begin{equation}\label{eq:observable-norm}
 \ip{\mathcal O(f)\Omega}{\mathcal O(g)\Omega},
 \qquad \norm{\mathcal O(f)\Omega}^2\geq0,
\end{equation}
provided these vectors exist in the positive physical Hilbert space.
Writing this formally as an expectation of $\mathcal O(f)^\dagger
\mathcal O(g)$ is justified only on its operator or quadratic-form
domain. Linearity of the state map supplies quadratic dependence on $f$;
the interaction strength of the theory does not affect this argument.

A Euclidean formulation may use a preparation $A_L(f)$ at positive
Euclidean time and reflection $\Theta$:
\begin{equation}\label{eq:os}
 B_{M,L}(f,g)=\langle(\Theta A_L(f))A_L(g)\rangle_M.
\end{equation}
Reflection positivity supplies its positive norm interpretation; a full
relativistic reconstruction requires the other Osterwalder--Schrader
hypotheses \cite{OS}. The arithmetic variable $x$ in
\eqref{eq:target} is not automatically Euclidean time.

An alternative classical construction is a nonnegative joint energy
$\mathcal E_M(\Phi)$ with boundary trace $f$, for which
$Q_L[f]=\inf_{\operatorname{Tr}\Phi=f}\mathcal E_M(\Phi)$.
Positivity only in the auxiliary Hessian is insufficient:
$w^2-2fw$ is coercive in $w$ and has minimum $-f^2$.
For nonlinear energies, the assertion that the entire minimized
boundary dependence is quadratic is itself a matching requirement.

\begin{proposition}[Conditional construction-to-RH implication]
\label{prop:conditional}
If the construction hypotheses for a specified model imply the positive
Hilbert space and domains in Definition~\ref{def:realization}, and the
complete matching identity is derived under those hypotheses, then
construction of that model with those observables implies RH.
It is also sufficient to have independently positive regulated pairings
$B_{\Lambda,L}$ such that, for every fixed pair of test functions and $L$,
\begin{equation}\label{eq:regulated-limit}
 \lim_{\Lambda\to\infty}B_{\Lambda,L}(f,g)=Q_L(f,g).
\end{equation}
\end{proposition}
\begin{proof}
The first assertion follows by taking $g=f$ in \eqref{eq:exact-match}
and applying Theorem~\ref{thm:weil}. In the second assertion the finite
limit of nonnegative real numbers $B_{\Lambda,L}(f,f)$ is nonnegative.
The same criterion applies. No rate uniform in all inputs is required.
\end{proof}

The proposition is a logical consequence of the classical criterion,
not the missing reduction theorem. We have not specified a model for
which its full premise is verified. An additive subtraction from a
regulated norm need not remain positive. Likewise, defining a covariance
to equal $Q_L$, or setting $\Phi=Q_L^{1/2}$ before positivity is known,
would reinsert the original problem into the construction hypothesis.

A common source rule compatible under enlargement of $I_L$ is not a
convenience. Separate positive realizations for every $L$ would suffice
logically for RH, but a realization at one fixed interval is not by itself
informative: as soon as $Q_L\geq0$ holds, $\Phi(f)=Q_L^{1/2}f$ into
$L^2(I_L)$ already satisfies Definition~\ref{def:realization} on that
interval. The content of a proposed realization is carried entirely by the
uniformity of its source rule in $L$, and by
Corollary~\ref{cor:spectral} that uniform object is the zero measure.
Short intervals are therefore places to test and discard a specified
candidate, not places where a partial result accumulates.

\begin{remark}[Stationary data at one interval]\label{rem:krein}
The same conclusion holds inside the narrower class of translation-covariant
sources. If $Q_L\geq0$ and $\varphi\in C_c^\infty$ has small support, then
$r\mapsto Q_L(\varphi,U_r\varphi)$ is a continuous positive definite
function on an interval, because
$\sum_{j,k}\bar c_jc_kQ_L(U_{r_j}\varphi,U_{r_k}\varphi)
=Q_L[\sum_kc_kU_{r_k}\varphi]\geq0$ whenever the translates stay inside
$I_L$. Krein's extension theorem \cite{Krein,Sasvari} then supplies a finite
positive measure whose Fourier transform agrees with it there, hence
stationary positive data on that scale. Positivity on one interval already
yields such data; only its compatibility as $L$ grows carries arithmetic
information. This remark is not used in any proof below.
\end{remark}

A mass gap is not needed for the elementary norm implication.
It must be imposed only if a chosen construction or observable limit
requires it. An extension of pure Yang--Mills by arithmetic defects or
extra fields would create construction obligations beyond the pure
Yang--Mills existence problem \cite{ClayYM}.

\subsection{Protected pairings and their adjoints}
Sphere and Schur quantization motivate Hilbert completions of protected
algebras. The basic data can be described as a complex algebra
$\mathcal A$, an antilinear automorphism $\rho$, and a twisted trace $T$
with positive Hermitian form
\begin{equation}\label{eq:twisted}
 H(a,b)=T(\rho(a)b),\qquad
 T(ab)=T(\rho^2(b)a).
\end{equation}
One quotients null vectors when necessary before completion. On the
algebraic core the right and left actions satisfy
$R_a^\dagger=L_{\rho(a)}$, since
$H(ba,c)=T(\rho(b)\rho(a)c)=H(b,\rho(a)c)$.
This does not identify $\rho$ with an arbitrary ordinary adjoint on a
single copy of the algebra, nor establish every unbounded closure.
The representation formulas used below come from
\cite{Sphere,Schur,RG}. Classification results for related positive
forms give independent mathematical support in specified algebraic
families \cite{EKRS,Klyuev}; they do not construct every correlator of
an interacting continuum field theory.

The intended conditional program allows an open but credible construction
problem. Explicit positive representations, stable regulated models,
protected localization formulas and compatible dual descriptions provide
different kinds of evidence. The source domain and matching identity
still have to be stated for the particular model. No general Schur
construction conjecture is promoted to a theorem in the calculations below.

\subsection{A necessary source-regularity condition}
\begin{proposition}[Finite-norm orbit smearing cannot match]\label{prop:bounded-source}
Fix $L>0$. A source of the form
$\Phi(f)=\int_{I_L}f(x)V_x\Omega\dd x$, with $V_x$ unitary,
$\Omega\in\HH$, and a well-defined strongly measurable orbit, cannot
have norm $Q_L[f]$ for all smooth test functions. The same holds for
a finite direct sum of such sources.
\end{proposition}
\begin{proof}
The estimate $\norm{\Phi(f)}\leq\sqrt L\norm\Omega\norm f$ makes
the proposed pairing bounded in the input $L^2$ norm. Let
$0\ne\chi\in C_c^\infty(I_L)$ and $f_N(x)=\chi(x)e^{iNx}$.
Equation~\eqref{eq:gamma-growth} and rapid decay of $\widehat\chi$
give
\begin{equation}\label{eq:high-frequency}
 \frac{Q_L[f_N]}{\log N}\longrightarrow\norm\chi^2.
\end{equation}
For completeness, after the change $\tau=N+u$ the integrand of the
gamma term divided by $\log N$ converges pointwise to
$|\widehat\chi(u)|^2$. The bound
$b((N+u)^2)\leq C[1+\log(2+N)+\log(2+|u|)]$ for $N\geq2$
permits dominated convergence. All other terms in
\eqref{eq:target} are bounded uniformly in $N$ on this interval.
This contradicts boundedness of the proposed norm.
\end{proof}
An exact source must be unbounded relative to the input $L^2$ norm.
An infinite family of insertions or a distributional field with finite
output on smooth tests is compatible with this requirement. The
proposition makes no assertion of global closability on the union of
all intervals, and imposes no obstruction on QFT as such.
